\section{Content-verified retrieval}
\label{sec:retrieval}
\label{sec:lanes}

Retrieval is the query module of the management system: it reads the
libraries, catalog and indices that acquisition and organization maintain, and
it returns either evidence for an answer or an explicit absence report.
Search, Librarian and Hybrid lanes use different retrieval affordances but
share a server-side content-verification tool when the calling workflow invokes
it. The tool fetches referenced document text, segments it and applies a local,
query-conditioned Laya verdict on a 0--1 scale---the model card's typed
binary-output field (\texttt{noul}), the calibrated probability that a
segment answers the posed evidence question; segments meeting the configured
threshold (0.5 by default) are returned as survivors with provenance. A scored segment below the
threshold is not admitted, but tool or fetch failures are reported separately
and a partially completed call may still return successfully scored survivors.

\subsection{Gate contract and retrieval lanes}
\label{sec:trace}
\label{sec:gate}
\label{sec:gate-primitive}
\label{sec:twostage}

The chat endpoint supplies two built-in policies. Lane A uses a wide dense search
(top 30, cosine threshold 0.35) within a pinned library and may escalate once
to the catalog ladder when no verified evidence survives. Lane B follows catalog descriptions, then
uses a deterministic scanner for literal or enumeration questions and the judge
for semantic questions. Endpoint-level timeouts and turn limits bound execution;
the lane call budgets and six-reference batching below are workflow instructions,
not independent server-enforced quotas. Exhaustion is intended to produce a
scope and blind-spot report. Lane C is a composed configuration used in
evaluation, not a third built-in server policy: it appends a hybrid directive
to the chat prompt that can call vector and catalog tools, merge up to twelve
candidates, judge them in batches of six, and return up to five survivors with
bounded evidence. Table~\ref{tab:lanes} and Fig.~\ref{fig:lanes} summarize the
policies and answer workflow.

The gate is the local Laya evidence classifier \cite{laya}, exposed as
\path{kb_laya_judge}. In reference mode, the server fetches each document,
judges its segments and returns provenance with a bounded evidence pack. The
verdict controls admission, not truth; the verifier remains subject to
calibration and transport limitations.

\begin{table*}[t]
\centering
\small
\setlength{\tabcolsep}{5pt}
\caption{Retrieval lanes as operating policies over a shared candidate and verification tool layer. Budgets describe the configured interface; latency and effectiveness are evaluated separately in Section~\ref{sec:eval}.}
\label{tab:lanes}
\begin{tabular}{@{}llll@{}}
\toprule
Lane & Retrieval affordance & Budget type & Intended terminal report \\
\midrule
Search & Pinned-library dense search & Endpoint turn limit & Answer or report not-found \\
Librarian & Catalog navigation and query-shaped scan/judging & Workflow call budget & Scope and blind spots \\
Hybrid & Agent-composed vector and catalog calls & Evaluation-specific candidate cap & Survivors for answer workflow \\
\bottomrule
\end{tabular}
\end{table*}

The two-stage tool builds a candidate pool with lexical search and optional
graph expansion, then applies dense ranking. It does not call the
content gate; a ranked chunk is not admissible evidence until verified, and a
similarity threshold is not an abstention mechanism. The 13 October catalog
incident showed that oversized tool returns can conceal documents before
reasoning begins. The tested repair adds a server-side character budget,
grouped views and pagination; its behavior is reported in the Evaluation
section.

\begin{figure*}[t]
\centering
\includegraphics[width=\textwidth]{fig2_lanes}
\caption{The three retrieval workflows can invoke the same content-verification
tool by different routes: a pinned wide dense search that may escalate once
to the catalog ladder when no survivors remain, a catalog ladder that selects a literal scanner or the judge,
and an agent-composed vector-and-catalog configuration capped at twelve
candidates in this evaluation. When invoked, the gate reads fetched text and
returns scored survivors with provenance; zero survivors or an exhausted
workflow budget is intended to elicit an explicit not-found report. Endpoint
tool and turn controls are server-side; the lane-level call sequence and
budgets are workflow instructions.}
\label{fig:lanes}
\end{figure*}

\subsection{Honest abstention}
A candidate list may contain high-similarity but unsupported documents. The
answer workflow instructs the agent to report not-found when its completed
evidence review finds no support, including the searched scope and known blind
spots, rather than falling back to parametric knowledge. This behavior was
observed in the reported trials; it is not a server-side semantic-classifier
guarantee. Timeouts, unavailable tools and incomplete searches are execution
failures, not evidence of semantic absence (Section~\ref{sec:eval}).
